\begin{titlepage}
    \centering
    \vspace*{2.5cm}

    {\zihao{1}\bfseries \ReportTitle \par}
    \vspace{1.2cm}
    {\zihao{2}\bfseries 项目名称：\ProjectName \par}
    \vspace{0.8cm}
    {\zihao{4}版本：\ReportVersion \qquad 日期：\ReportDate \par}

    \vspace{2.0cm}

    \renewcommand{\arraystretch}{1.5}
    \begin{tabular}{>{\raggedleft\arraybackslash}p{0.24\textwidth} p{0.56\textwidth}}
        参赛学校/单位： & \SchoolName \\
        团队名称： & \TeamName \\
        成员信息： & \MemberList \\
        文档定位： & 初赛提交版设计说明书 \\
        适用范围： & 设计方案、技术路线、验证证据、FPGA pre-board 状态 \\
    \end{tabular}

    \vfill

    {\zihao{5}
    本文档围绕 \texttt{YH\_rv\_cpu} 当前比赛主线工程编写，目标是在初赛阶段提供一份
    可直接提交、可继续扩展、且和仓库真实状态一致的设计说明书。文中区分了冻结比赛基线、
    活跃扩展验证和外部阻塞，避免将实验态结果误写为正式结论。 \par}
\end{titlepage}

\section*{摘\quad 要}
\addcontentsline{toc}{section}{摘要}

\texttt{YH\_rv\_cpu} 是一套面向比赛提交的工程化 RISC-V CPU 基线，而不是一次性实验核。
当前主线以 \texttt{RV32I + Zicsr} 为冻结比赛口径，采用五级流水结构，配套 SoC ROM/RAM、
UART、定时器与 DONE 标志，形成了可以直接运行 Smoke、\texttt{riscv-tests}、CoreMark、
Vivado \texttt{impl50} 以及 FPGA-like probe 的完整工程闭环。为了提升初赛提交材料的完整性，
本说明书在不夸大设计边界的前提下，系统整理了项目定位、微架构设计、验证矩阵、性能结果、
FPGA pre-board 状态和后续优化规划。

截至 2026 年 4 月，项目已经完成 frozen baseline 与 fresh evidence 的双层收口：CoreMark short
达到 \texttt{0.912472 CoreMark/MHz}，strict \texttt{>=10s} 路径达到
\texttt{0.912465 CoreMark/MHz}；\texttt{rv32 baseline} 与 \texttt{rv64 baseline}
分别通过 \texttt{33/33} 和 \texttt{21/21}；扩展 \texttt{full-ui} 验证达到
\texttt{rv32 42/42} 与 \texttt{rv64 54/54}；Vivado \texttt{impl50} 结果为
\texttt{2556 LUT / 2170 FF / 4 BRAM / 0 DSP}，并保留正时序裕量。与此同时，项目也明确记录了
当前未完成项主要集中在实体板卡和最终板级 I/O delay 约束，属于外部阻塞，而不是初赛文档交付阻塞。

本说明书的价值不只在于展示结果，更在于解释设计取舍与工程方法：哪些优化被保留，
哪些方向被证据否定，为什么当前冻结口径仍维持 \texttt{RV32I + Zicsr}，以及为什么初赛阶段
应以设计/技术文档为核心交付、以性能与验证材料作为增强支撑。

\section*{关键词}
\addcontentsline{toc}{section}{关键词}

RISC-V；五级流水；CoreMark；FPGA 原型；比赛提交；工程化验证
\clearpage
